\section{A tensor estimate for the divergence tests}\label{sec:tensor}

The gradient and Hessian pairings in \eqref{eq:test-pairing} will be
estimated using two quadratic energies, one for $W$ and one for $B$.
The following lemma bounds their sum by $\E u^2$. Its proof uses weighted
integration by parts and the moment equation; homogeneity enters only
when the lemma is applied to the summands $\psi_i$ in
Section~\ref{sec:variance}.

\begin{lemma}[Moment-coordinate tensor estimate]\label{lem:tensor}
Let $d\mu=Z^{-1}e^{-V}dx$ and $d\nu=e^{-\varphi}dz$ be smooth
probability measures on $\R^n$.  Suppose $D_z^2\varphi$ is positive
definite and $x=\nabla\varphi(z)$ is a smooth diffeomorphism pushing
$\nu$ to $\mu$.
For a real $f\in C_c^\infty(\R_x^n)$, put
\[
 \tau=D_z^2\varphi,\quad M=\tau^{-1},\quad H=D_x^2V,
 \quad B=\tau(D_x^2f)\tau,
 \quad W=\delta_\mu B,
 \quad u=\delta_\mu W.
\]
View source fields in target coordinates by the inverse moment map.
Then
\begin{equation}\label{eq:tensor-estimate}
 \E\tr\bigl((D_xW)^2\bigr)
 \ge \E\tr(MBHB),
\end{equation}
and
\begin{equation}\label{eq:tensor-energy}
 \E u^2\ge\E\tr(MBHB)+\E W^THW.
\end{equation}
Here $(D_xW)_{ij}=\partial_{x_j}W_i$, and the square in
\eqref{eq:tensor-estimate} is a matrix square.
\end{lemma}

\begin{proof}
All test fields used in the integrations by parts below have compact support.
Weighted Bochner calculations underlie many variance estimates; see
Bakry--\'Emery \cite[Proposition~3]{BakryEmery1985} and the
transport-Hessian formulas in \cite[Section~4.3]{KLR2026}.
Here the required identity is for a vector field, so we derive it
explicitly and retain the matrix trace $\tr((D_xW)^2)$.
First, for any such vector field $W$ and $u=\delta_\mu W$,
\[
 \partial_{x_i}u
 =H_{ij}W_j+(V_{x_j}-\partial_{x_j})\partial_{x_i}W_j.
\]
Pairing with $W_i$ and using adjunction gives the weighted identity
\begin{equation}\label{eq:bochner}
 \E u^2
 =\E W_i\partial_{x_i}u
 =\E W^THW
   +\E(\partial_{x_j}W_i)(\partial_{x_i}W_j)
 =\E W^THW+\E\tr\bigl((D_xW)^2\bigr).
\end{equation}
Thus it remains to prove \eqref{eq:tensor-estimate}.
Its proof compares two expansions, integrates their derivative terms
by parts, and writes the resulting difference as a sum of squares.

Recall $h=D_x^2f$ and $P=\tau h$ from \eqref{eq:test-fields}, and let
$(C_i)_{kl}=C_{ikl}$, where $C=D_z^3\varphi$.
The moment equation \eqref{eq:moment} reads
\[
 V(x(z))=\varphi(z)+\log\det\tau(z)-\log Z.
\]
Its first derivative gives
\[
 \tau_{ik}V_{x_k}=x_i+c_i,
 \qquad c_i=\tr(MC_i).
\]
Since $\partial_jM=-MC_jM$, differentiating once more gives
\[
 (\tau H\tau)_{ij}
 =\tau_{ij}+M_{kl}\varphi_{klij}
   -\tr(MC_iMC_j)-C_{ijk}V_{x_k}.
\]
To combine the fourth-derivative term and the term involving
$\nabla_xV$ into a weighted divergence, note that
\[
 \partial_dM_{dk}
 =-M_{da}C_{abd}M_{bk}=-c_bM_{bk},
 \qquad
 \partial_dM_{dk}-x_dM_{dk}=-V_{x_k}.
\]
Full symmetry of $D_z^4\varphi$ therefore yields
\begin{equation}\label{eq:twice-moment}
 (\tau H\tau-\tau)_{ij}
 =(\partial_d-x_d)(M_{dk}C_{ijk})
   -\tr(MC_iMC_j).
\end{equation}

For the other expansion, the commutator
$[\partial_l,\partial_k^*]=\tau_{kl}$ gives
\[
 D_zW=B+R,\qquad R_{al}=\partial_k^*\partial_lP_{ak}.
\]
Both $D_zW$ and $B$ are symmetric, so $R$ is symmetric.
Since $D_xW=(B+R)M$ and $MBM=h$, expansion followed by one
integration by parts gives
\begin{equation}\label{eq:trace-expansion}
 \E\tr\bigl((D_xW)^2\bigr)
 =\E\tr(MBMB)
  +2\E(\partial_kh_{al})(\partial_lP_{ak})
  +\E\tr(MRMR).
\end{equation}
The last term is nonnegative:
$\tr(MRMR)=\|M^{1/2}RM^{1/2}\|_{\mathrm{HS}}^2$.
On the right side of \eqref{eq:tensor-estimate}, the coefficient of
$\tau H\tau$ in the trace contraction is
\[
 Q=MBMBM=h\tau h.
\]
Substituting \eqref{eq:twice-moment} and integrating its divergence
term by parts gives
\begin{equation}\label{eq:H-expansion}
 \E\tr(MBHB)
 =\E\tr(MBMB)
   -\E(\partial_dQ_{ij})M_{dk}C_{ijk}
   -\E Q_{ij}\tr(MC_iMC_j).
\end{equation}
The common first term cancels.  We now evaluate the three remaining
contracted terms in coordinates where $\tau=I$ at the point under
consideration.

This normalization uses a constant dual change of coordinates.
For an invertible constant matrix $T$, set
\[
 z=Tz',\qquad x'=T^Tx,\qquad
 f'(x')=f(T^{-T}x').
\]
The corresponding normalized potentials and partition function are
\[
 \varphi'(z')=\varphi(Tz')-\log|\det T|,
 \qquad V'(x')=V(T^{-T}x'),
 \qquad Z'=|\det T|Z.
\]
They preserve both probability measures by change of variables,
$x'=\nabla_{z'}\varphi'$, and the moment equation.  The chain rule gives
\[
 \begin{array}{lll}
 \tau'=T^T\tau T,& B'=T^TBT,& R'=T^TRT,\\
 M'=T^{-1}MT^{-T},& h'=T^{-1}hT^{-T},& H'=T^{-1}HT^{-T},\\
 Q'=T^{-1}QT^{-T},& P'=T^TPT^{-T},& W'=T^TW.
 \end{array}
\]
Each $z$-derivative index, including all three indices of $C$,
transforms covariantly with $T$; the same law holds for $\partial_k^*$.
Also $D_{x'}W'=T^T(D_xW)T^{-T}$.
It follows that every contracted scalar in
\eqref{eq:trace-expansion} and \eqref{eq:H-expansion} is invariant.
For example, in
$(\partial_kh_{al})(\partial_lP_{ak})$ each index pairs one
covariant and one contravariant factor.
At a fixed point we may therefore take $T=\tau^{-1/2}$, using
the value of $\tau$ at that point.  All derivatives are computed
under this constant transformation.  No moving frame is
differentiated; the integrations by parts have already been completed.

In these normalized coordinates, put
\[
 S_{ijk}=f_{x_ix_jx_k},\qquad J_{ijk}=h_{ia}C_{ajk}.
\]
The tensor $S$ is fully symmetric, and $J$ is symmetric in its last
two indices.  Write $J^{(12)}_{ijk}=J_{jik}$, and use the Euclidean
tensor norm and inner product in this frame.
The three contractions are
\begin{align}
 2(\partial_kh_{al})(\partial_lP_{ak})
   &=2\|S\|^2+2\langle S,J\rangle,\label{eq:normalized-cross}\\
 (\partial_dQ_{ij})C_{ijd}
   &=2\langle S,J\rangle+\langle J,J^{(12)}\rangle,
       \label{eq:normalized-derivative}\\
 Q_{ij}\tr(C_iC_j)&=\|J\|^2.\label{eq:normalized-cubic}
\end{align}
For \eqref{eq:normalized-cross}, use
$\partial_kh_{al}=S_{alk}$ and
$\partial_lP_{ak}=C_{ail}h_{ik}+S_{akl}$.
For \eqref{eq:normalized-derivative}, differentiating the two factors
$h$ in $Q=h\tau h$ gives the two copies of $\langle S,J\rangle$.
Differentiating its middle factor gives
\[
 \sum_{iabjd}h_{ia}C_{abd}h_{bj}C_{ijd}
 =\sum_d\tr(hC_dhC_d)
 =\langle J,J^{(12)}\rangle.
\]
Finally, \eqref{eq:normalized-cubic} follows by expanding
$Q=h^2$ and the trace, using the full symmetry of $C$.

Thus, apart from the nonnegative $R$ term, the integrand of
\eqref{eq:trace-expansion} minus \eqref{eq:H-expansion} is
\[
 2\|S\|^2+4\langle S,J\rangle
   +\|J\|^2+\langle J,J^{(12)}\rangle.
\]
Full symmetrization of $J$ is
\[
 (\operatorname{Sym}J)_{ijk}
 =\frac{J_{ijk}+J_{jik}+J_{kij}}3.
\]
Its norm satisfies
\[
 \|\operatorname{Sym}J\|^2
 =\frac{\|J\|^2+2\langle J,J^{(12)}\rangle}{3},
 \qquad
 \langle S,\operatorname{Sym}J\rangle=\langle S,J\rangle.
\]
Consequently the preceding integrand equals
\begin{equation}\label{eq:tensor-squares}
 2\|S+\operatorname{Sym}J\|^2
 +\frac16\|J-J^{(12)}\|^2\ge0.
\end{equation}
This pointwise calculation proves \eqref{eq:tensor-estimate};
combining it with \eqref{eq:bochner} proves
\eqref{eq:tensor-energy}.
\end{proof}

For the soft slab potentials, $H$ is the sum of a positive quadratic
part and the positive semidefinite Hessians of the homogeneous
summands.  Lemma~\ref{lem:tensor} will control each summand through
the two terms on the right side of \eqref{eq:tensor-energy}.
